% =========================================================
% PHẦN A. TRẮC NGHIỆM 4 LỰA CHỌN
% =========================================================

% ID: L12C1B1-001-TN
% Mức: NB
%=== Câu 1 ===%
\dangbt{Phân loại chất rắn}
\begin{ex}
Chất rắn kết tinh là chất rắn
\choice
{không có cấu trúc tinh thể}
{\True có cấu trúc tinh thể xác định}
{luôn có nhiệt độ nóng chảy không xác định}
{luôn mềm dần khi được đun nóng}
\loigiai{
Chất rắn kết tinh có cấu trúc tinh thể với sự sắp xếp các hạt có trật tự và tuần hoàn trong không gian.
}
\end{ex}

% ID: L12C1B1-002-TN
% Mức: NB
%=== Câu 2 ===%
\dangbt{Phân loại chất rắn}
\begin{ex}
Chất nào sau đây là chất rắn kết tinh?
\choice
{\True Muối ăn}
{Thủy tinh}
{Nhựa đường}
{Cao su}
\loigiai{
Muối ăn NaCl có cấu trúc tinh thể xác định nên là chất rắn kết tinh. Thủy tinh, nhựa đường và cao su là những chất rắn vô định hình.
}
\end{ex}

% ID: L12C1B1-003-TN
% Mức: NB
%=== Câu 3 ===%
\dangbt{Tính chất chất rắn kết tinh}
\begin{ex}
Đặc điểm nào sau đây đúng với chất rắn kết tinh?
\choice
{Không có cấu trúc tinh thể}
{\True Có nhiệt độ nóng chảy xác định}
{Luôn mềm dần khi nóng lên}
{Không có trật tự trong sự sắp xếp các hạt}
\loigiai{
Chất rắn kết tinh có cấu trúc tinh thể xác định và có nhiệt độ nóng chảy xác định.
}
\end{ex}

% ID: L12C1B1-004-TN
% Mức: NB
%=== Câu 4 ===%
\dangbt{Tính chất chất rắn vô định hình}
\begin{ex}
Đặc điểm nào sau đây là đặc trưng của chất rắn vô định hình?
\choice
{Có cấu trúc tinh thể xác định}
{Có nhiệt độ nóng chảy xác định}
{\True Không có nhiệt độ nóng chảy xác định}
{Các hạt sắp xếp hoàn toàn tuần hoàn}
\loigiai{
Chất rắn vô định hình không có cấu trúc tinh thể xác định và không có nhiệt độ nóng chảy xác định. Khi đun nóng, chúng thường mềm dần trong một khoảng nhiệt độ.
}
\end{ex}

% ID: L12C1B1-005-TN
% Mức: NB
%=== Câu 5 ===%
\dangbt{Ví dụ về chất rắn}
\begin{ex}
Chất nào sau đây thường được xem là chất rắn vô định hình?
\choice
{Kim cương}
{Muối ăn}
{Đồng}
{\True Thủy tinh}
\loigiai{
Thủy tinh là chất rắn vô định hình. Kim cương, muối ăn và đồng là các chất rắn kết tinh.
}
\end{ex}

% ID: L12C1B1-006-TN
% Mức: NB
%=== Câu 6 ===%
\dangbt{Nhiệt độ nóng chảy}
\begin{ex}
Khi đun nóng một chất rắn kết tinh, chất rắn chuyển sang trạng thái lỏng
\choice
{trong một khoảng nhiệt độ rất rộng}
{\True tại một nhiệt độ xác định ở điều kiện xác định}
{ở mọi nhiệt độ}
{không phụ thuộc vào áp suất}
\loigiai{
Ở điều kiện xác định, chất rắn kết tinh nóng chảy tại một nhiệt độ xác định gọi là nhiệt độ nóng chảy.
}
\end{ex}

% ID: L12C1B1-007-TN
% Mức: NB
%=== Câu 7 ===%
\dangbt{Cấu trúc chất rắn}
\begin{ex}
Trong chất rắn kết tinh, các hạt cấu tạo nên chất
\choice
{chuyển động hỗn loạn hoàn toàn}
{\True được sắp xếp có trật tự trong mạng tinh thể}
{không có tương tác với nhau}
{luôn đứng yên tuyệt đối}
\loigiai{
Các hạt trong chất rắn kết tinh được sắp xếp có trật tự và tạo thành mạng tinh thể.
}
\end{ex}

% ID: L12C1B1-008-TN
% Mức: NB
%=== Câu 8 ===%
\dangbt{Phân biệt chất rắn}
\begin{ex}
Phát biểu nào sau đây đúng?
\choice
{Mọi chất rắn đều có nhiệt độ nóng chảy xác định}
{\True Chất rắn kết tinh có nhiệt độ nóng chảy xác định}
{Chất rắn vô định hình có cấu trúc tinh thể hoàn chỉnh}
{Chất rắn vô định hình nóng chảy đột ngột tại một nhiệt độ xác định}
\loigiai{
Chất rắn kết tinh có nhiệt độ nóng chảy xác định, còn chất rắn vô định hình thường mềm dần khi được đun nóng.
}
\end{ex}

% ID: L12C1B1-009-TN
% Mức: TH
%=== Câu 9 ===%
\dangbt{So sánh chất rắn}
\begin{ex}
Cho các chất: kim cương, thủy tinh, nhôm, cao su. Số chất rắn kết tinh là
\choice
{$1$}
{$2$}
{\True $2$}
{$4$}
\loigiai{
Kim cương và nhôm là chất rắn kết tinh; thủy tinh và cao su là chất rắn vô định hình. Vậy có $2$ chất rắn kết tinh.
\end{ex}

% ID: L12C1B1-010-TN
% Mức: TH
%=== Câu 10 ===%
\dangbt{Nhận biết chất rắn}
\begin{ex}
Một chất rắn khi được đun nóng không chuyển sang trạng thái lỏng tại một nhiệt độ xác định mà mềm dần trong một khoảng nhiệt độ. Chất đó có khả năng là
\choice
{kim cương}
{đồng}
{\True thủy tinh}
{muối ăn}
\loigiai{
Hiện tượng mềm dần trong một khoảng nhiệt độ là đặc trưng của chất rắn vô định hình. Thủy tinh là chất rắn vô định hình.
}
\end{ex}

% ID: L12C1B1-011-TN
% Mức: TH
%=== Câu 11 ===%
\dangbt{Phân loại chất rắn}
\begin{ex}
Trong các chất sau: sắt, thủy tinh, nhựa đường, muối ăn, cao su. Có bao nhiêu chất rắn kết tinh?
\choice
{$1$}
{\True $2$}
{$3$}
{$4$}
\loigiai{
Sắt và muối ăn là chất rắn kết tinh. Thủy tinh, nhựa đường và cao su là chất rắn vô định hình. Vậy có $2$ chất.
}
\end{ex}

% ID: L12C1B1-012-TN
% Mức: TH
%=== Câu 12 ===%
\dangbt{Tính chất chất rắn vô định hình}
\begin{ex}
Khi đun nóng thủy tinh, hiện tượng nào sau đây xảy ra?
\choice
{Thủy tinh nóng chảy hoàn toàn ngay tại một nhiệt độ xác định}
{\True Thủy tinh mềm dần trong một khoảng nhiệt độ}
{Thủy tinh không thay đổi trạng thái}
{Thủy tinh luôn giữ nguyên độ cứng}
\loigiai{
Thủy tinh là chất rắn vô định hình nên khi đun nóng thường mềm dần trong một khoảng nhiệt độ.
}
\end{ex}

% ID: L12C1B1-013-TN
% Mức: TH
%=== Câu 13 ===%
\dangbt{So sánh chất rắn}
\begin{ex}
Điểm khác biệt cơ bản giữa chất rắn kết tinh và chất rắn vô định hình là
\choice
{cả hai đều có cấu trúc tinh thể}
{chất rắn kết tinh không có nhiệt độ nóng chảy}
{\True chất rắn kết tinh có cấu trúc tinh thể và nhiệt độ nóng chảy xác định}
{chất rắn vô định hình luôn cứng hơn chất rắn kết tinh}
\loigiai{
Chất rắn kết tinh có cấu trúc tinh thể xác định và nhiệt độ nóng chảy xác định; chất rắn vô định hình không có cấu trúc tinh thể xác định và không có nhiệt độ nóng chảy xác định.
}
\end{ex}

% ID: L12C1B1-014-TN
% Mức: TH
%=== Câu 14 ===%
\dangbt{Ứng dụng chất rắn vô định hình}
\begin{ex}
Thủy tinh được sử dụng để chế tạo nhiều sản phẩm có hình dạng phức tạp chủ yếu vì
\choice
{thủy tinh luôn ở trạng thái lỏng}
{\True khi nung nóng, thủy tinh mềm dần và có thể tạo hình}
{thủy tinh có cấu trúc tinh thể hoàn hảo}
{thủy tinh có nhiệt độ nóng chảy duy nhất}
\loigiai{
Thủy tinh là chất rắn vô định hình. Khi nung nóng, thủy tinh mềm dần trong một khoảng nhiệt độ nên thuận lợi cho việc tạo hình.
}
\end{ex}

% ID: L12C1B1-015-TN
% Mức: VD
%=== Câu 15 ===%
\dangbt{Giải thích tính chất chất rắn}
\begin{ex}
Một người nung nóng hai mẫu vật A và B. Mẫu A chuyển đột ngột từ rắn sang lỏng tại $T_1$, còn mẫu B mềm dần trong khoảng nhiệt độ từ $T_2$ đến $T_3$. Kết luận nào sau đây phù hợp?
\choice
{Cả A và B đều là chất rắn kết tinh}
{Cả A và B đều là chất rắn vô định hình}
{\True A có thể là chất rắn kết tinh, B có thể là chất rắn vô định hình}
{A có thể là chất rắn vô định hình, B có thể là chất rắn kết tinh}
\loigiai{
Mẫu A nóng chảy tại một nhiệt độ xác định nên có thể là chất rắn kết tinh. Mẫu B mềm dần trong một khoảng nhiệt độ nên có thể là chất rắn vô định hình.
}
\end{ex}

% ID: L12C1B1-016-TN
% Mức: VD
%=== Câu 16 ===%
\dangbt{Nhiệt độ nóng chảy}
\begin{ex}
Hai chất rắn kết tinh khác nhau có thể có cùng nhiệt độ nóng chảy hay không?
\choice
{Không, vì mỗi chất rắn kết tinh luôn có nhiệt độ nóng chảy khác nhau}
{\True Có, vì các chất khác nhau có thể có cùng một giá trị nhiệt độ nóng chảy}
{Không, vì nhiệt độ nóng chảy không phụ thuộc bản chất chất}
{Có, nhưng chỉ khi chúng là cùng một chất}
\loigiai{
Mỗi chất rắn kết tinh có nhiệt độ nóng chảy đặc trưng ở một điều kiện xác định, nhưng hai chất khác nhau về nguyên tắc vẫn có thể có cùng giá trị nhiệt độ nóng chảy.
}
\end{ex}

% ID: L12C1B1-017-TN
% Mức: VD
%=== Câu 17 ===%
\dangbt{Ứng dụng tính chất chất rắn}
\begin{ex}
Một cơ sở sản xuất muốn tạo hình một vật bằng thủy tinh. Phương pháp phù hợp nhất là
\choice
{làm lạnh thủy tinh đến nhiệt độ rất thấp rồi uốn}
{\True nung nóng để thủy tinh mềm dần rồi tạo hình}
{nung thủy tinh đến một nhiệt độ xác định rồi mới tạo hình trong trạng thái rắn}
{không cần nung nóng}
\loigiai{
Thủy tinh là chất rắn vô định hình. Khi nung nóng, thủy tinh mềm dần nên có thể tạo hình thuận lợi.
}
\end{ex}

% ID: L12C1B1-018-TN
% Mức: VD
%=== Câu 18 ===%
\dangbt{Giải thích hiện tượng}
\begin{ex}
Một vật liệu được nung nóng và có thể uốn, nắn trong một khoảng nhiệt độ khá rộng trước khi chuyển hoàn toàn sang trạng thái lỏng. Vật liệu này có đặc điểm
\choice
{có mạng tinh thể hoàn chỉnh}
{có nhiệt độ nóng chảy xác định}
{\True không có nhiệt độ nóng chảy xác định}
{luôn là kim loại}
\loigiai{
Vật liệu mềm dần trong một khoảng nhiệt độ là đặc trưng của chất rắn vô định hình. Vì vậy, vật liệu không có nhiệt độ nóng chảy xác định.
}
\end{ex}

% =========================================================
% PHẦN B. ĐÚNG / SAI
% =========================================================

% ID: L12C1B1-019-DS
% Mức: NB
%=== Câu 19 ===%
\dangbt{Phân biệt chất rắn kết tinh và vô định hình}
\begin{ex}
Xét các phát biểu sau:
\choiceTF
{\True Chất rắn kết tinh có cấu trúc tinh thể xác định.}
{\True Chất rắn kết tinh có nhiệt độ nóng chảy xác định ở điều kiện xác định.}
{Chất rắn vô định hình có cấu trúc mạng tinh thể tuần hoàn xác định.}
{Chất rắn vô định hình luôn nóng chảy đột ngột tại một nhiệt độ xác định.}
\loigiai{
\begin{itemchoice}
\item (Đúng) Chất rắn kết tinh có cấu trúc tinh thể xác định.
\item (Đúng) Chất rắn kết tinh có nhiệt độ nóng chảy xác định ở điều kiện xác định.
\item (Sai) Chất rắn vô định hình không có cấu trúc mạng tinh thể tuần hoàn xác định.
\item (Sai) Chất rắn vô định hình thường mềm dần trong một khoảng nhiệt độ.
\end{itemchoice}
}
\end{ex}

% ID: L12C1B1-020-DS
% Mức: NB
%=== Câu 20 ===%
\dangbt{Nhận biết chất rắn}
\begin{ex}
Xét các chất: thủy tinh, kim cương, muối ăn và cao su.
\choiceTF
{\True Thủy tinh là chất rắn vô định hình.}
{\True Kim cương là chất rắn kết tinh.}
{\True Muối ăn là chất rắn kết tinh.}
{Cao su là chất rắn kết tinh có nhiệt độ nóng chảy xác định.}
\loigiai{
\begin{itemchoice}
\item (Đúng) Thủy tinh là chất rắn vô định hình.
\item (Đúng) Kim cương có cấu trúc tinh thể nên là chất rắn kết tinh.
\item (Đúng) Muối ăn NaCl có cấu trúc tinh thể.
\item (Sai) Cao su là chất rắn vô định hình.
\end{itemchoice}
}
\end{ex}

% ID: L12C1B1-021-DS
% Mức: TH
%=== Câu 21 ===%
\dangbt{Tính chất và ứng dụng}
\begin{ex}
Một người thợ sử dụng thủy tinh để chế tạo các sản phẩm có hình dạng phức tạp.
\choiceTF
{\True Thủy tinh có thể mềm dần khi được nung nóng.}
{\True Có thể tạo hình thủy tinh khi vật liệu đang ở trạng thái mềm thích hợp.}
{Thủy tinh phải có một nhiệt độ nóng chảy xác định như kim loại kết tinh.}
{\True Tính chất mềm dần khi nung nóng giúp thủy tinh thuận lợi trong quá trình tạo hình.}
\loigiai{
\begin{itemchoice}
\item (Đúng) Thủy tinh là chất rắn vô định hình và mềm dần khi nung nóng.
\item (Đúng) Khi đạt trạng thái mềm thích hợp, thủy tinh có thể được tạo hình.
\item (Sai) Thủy tinh không có nhiệt độ nóng chảy xác định.
\item (Đúng) Đây là một trong những tính chất thuận lợi của thủy tinh trong chế tạo sản phẩm.
\end{itemchoice}
}
\end{ex}

% ID: L12C1B1-022-DS
% Mức: TH
%=== Câu 22 ===%
\dangbt{So sánh chất rắn}
\begin{ex}
Xét các phát biểu sau:
\choiceTF
{\True Chất rắn kết tinh có sự sắp xếp các hạt có trật tự trong mạng tinh thể.}
{Chất rắn vô định hình luôn có cấu trúc hoàn toàn hỗn loạn ở mọi khoảng cách.}
{\True Chất rắn vô định hình không có nhiệt độ nóng chảy xác định.}
{\True Khi nung nóng, chất rắn vô định hình thường mềm dần trong một khoảng nhiệt độ.}
\loigiai{
\begin{itemchoice}
\item (Đúng) Các hạt trong chất rắn kết tinh được sắp xếp có trật tự.
\item (Sai) Không nên hiểu chất rắn vô định hình là hoàn toàn không có bất kỳ trật tự nào ở mọi khoảng cách; đặc trưng quan trọng là không có trật tự tinh thể dài hạn.
\item (Đúng) Chất rắn vô định hình không có nhiệt độ nóng chảy xác định.
\item (Đúng) Chúng thường mềm dần trong một khoảng nhiệt độ khi được nung nóng.
\end{itemchoice}
}
\end{ex}

% =========================================================
% PHẦN C. TRẢ LỜI NGẮN
% =========================================================

% ID: L12C1B1-023-TLN
% Mức: NB
%=== Câu 23 ===%
\dangbt{Phân loại chất rắn}
\begin{bt}
Trong các chất: kim cương, thủy tinh, đồng, cao su, muối ăn và nhựa đường, có bao nhiêu chất là chất rắn kết tinh?
\shortans[oly]{$3$}
\loigiai{
Kim cương, đồng và muối ăn là chất rắn kết tinh. Thủy tinh, cao su và nhựa đường là chất rắn vô định hình. Vậy có $3$ chất.
}
\end{bt}

% ID: L12C1B1-024-TLN
% Mức: NB
%=== Câu 24 ===%
\dangbt{Đặc điểm chất rắn}
\begin{bt}
Một chất rắn kết tinh có nhiệt độ nóng chảy xác định. Khi nung nóng, chất bắt đầu nóng chảy ở $120^\circ\mathrm{C}$. Nhiệt độ nóng chảy của chất là bao nhiêu độ Celsius?
\shortans[oly]{$120$}
\loigiai{
Nhiệt độ tại đó chất rắn kết tinh bắt đầu nóng chảy trong điều kiện đã cho chính là nhiệt độ nóng chảy. Do đó:

$$
t_{\mathrm{nc}}=120^\circ\mathrm{C}.
$$

}
\end{bt}

% ID: L12C1B1-025-TLN
% Mức: TH
%=== Câu 25 ===%
\dangbt{Phân loại chất rắn}
\begin{bt}
Cho các chất: sắt, thủy tinh, nhôm, cao su, muối ăn. Có bao nhiêu chất là chất rắn vô định hình?
\shortans[oly]{$2$}
\loigiai{
Thủy tinh và cao su là chất rắn vô định hình. Sắt, nhôm và muối ăn là chất rắn kết tinh. Vậy có $2$ chất.
}
\end{bt}

% ID: L12C1B1-026-TLN
% Mức: TH
%=== Câu 26 ===%
\dangbt{Phân biệt chất rắn}
\begin{bt}
Cho bốn phát biểu:

(1) Chất rắn kết tinh có cấu trúc tinh thể xác định.

(2) Chất rắn kết tinh có nhiệt độ nóng chảy xác định.

(3) Chất rắn vô định hình không có nhiệt độ nóng chảy xác định.

(4) Chất rắn vô định hình luôn có cấu trúc tinh thể.

Có bao nhiêu phát biểu đúng?
\shortans[oly]{$3$}
\loigiai{
Các phát biểu (1), (2), (3) đúng; phát biểu (4) sai. Vậy có $3$ phát biểu đúng.
}
\end{bt}

% ID: L12C1B1-027-TLN
% Mức: TH
%=== Câu 27 ===%
\dangbt{Nhận biết qua hiện tượng}
\begin{bt}
Một vật liệu khi nung nóng mềm dần trong khoảng từ $600^\circ\mathrm{C}$ đến $700^\circ\mathrm{C$}$ thay vì nóng chảy đột ngột tại một nhiệt độ xác định. Vật liệu đó thuộc loại chất rắn nào?
\shortans[oly]{vô định hình}
\loigiai{
Vật liệu mềm dần trong một khoảng nhiệt độ và không có nhiệt độ nóng chảy xác định nên đó là chất rắn vô định hình.
}
\end{bt}

% ID: L12C1B1-028-TLN
% Mức: VD
%=== Câu 28 ===%
\dangbt{Ứng dụng tính chất chất rắn}
\begin{bt}
Một người thợ có sáu mẫu vật liệu, trong đó có kim cương, thủy tinh, đồng, cao su, muối ăn và nhựa đường. Người thợ muốn chọn các vật liệu có thể mềm dần khi nung nóng để thuận lợi cho việc tạo hình. Có bao nhiêu mẫu vật liệu phù hợp?
\shortans[oly]{$3$}
\loigiai{
Các chất rắn vô định hình trong danh sách là thủy tinh, cao su và nhựa đường. Đây là các vật liệu có xu hướng mềm dần khi nung nóng. Vậy có $3$ mẫu phù hợp.
}
\end{bt}

% =========================================================
% PHẦN D. TỰ LUẬN
% =========================================================

% ID: L12C1B1-029-TL
% Mức: NB
%=== Câu 29 ===%
\dangbt{Khái niệm chất rắn kết tinh}
\begin{bt}
Nêu khái niệm chất rắn kết tinh và trình bày hai đặc điểm cơ bản của chất rắn kết tinh.
\loigiai{
Chất rắn kết tinh là chất rắn có cấu trúc tinh thể, trong đó các hạt cấu tạo nên chất được sắp xếp có trật tự và tuần hoàn trong không gian.

Hai đặc điểm cơ bản là:
\begin{itemize}
\item Có cấu trúc tinh thể xác định.
\item Có nhiệt độ nóng chảy xác định ở điều kiện xác định.
\end{itemize}
}
\end{bt}

% ID: L12C1B1-030-TL
% Mức: TH
%=== Câu 30 ===%
\dangbt{Phân biệt chất rắn kết tinh và vô định hình}
\begin{bt}
Phân biệt chất rắn kết tinh và chất rắn vô định hình về cấu trúc và đặc điểm nóng chảy. Cho mỗi loại một ví dụ.
\loigiai{
\begin{itemize}
\item \textbf{Chất rắn kết tinh:} có cấu trúc tinh thể xác định, các hạt được sắp xếp có trật tự. Chất rắn kết tinh có nhiệt độ nóng chảy xác định ở điều kiện xác định. Ví dụ: muối ăn, kim loại đồng, kim cương.
\item \textbf{Chất rắn vô định hình:} không có cấu trúc tinh thể xác định. Khi đun nóng, chất thường mềm dần trong một khoảng nhiệt độ và không có nhiệt độ nóng chảy xác định. Ví dụ: thủy tinh, cao su, nhựa đường.
\end{itemize}
}
\end{bt}

% ID: L12C1B1-031-TL
% Mức: TH
%=== Câu 31 ===%
\dangbt{Giải thích tính chất chất rắn}
\begin{bt}
Giải thích tại sao khi nung nóng thủy tinh, người thợ có thể uốn, nắn và tạo hình thủy tinh trong một khoảng nhiệt độ, trong khi đối với một chất rắn kết tinh như đồng thì hiện tượng này không giống như vậy.
\loigiai{
Thủy tinh là chất rắn vô định hình. Nó không có cấu trúc tinh thể tuần hoàn xác định và không có nhiệt độ nóng chảy xác định. Khi nung nóng, thủy tinh mềm dần trong một khoảng nhiệt độ, do đó người thợ có thể uốn, nắn và tạo hình.

Đồng là chất rắn kết tinh, có cấu trúc tinh thể xác định và nhiệt độ nóng chảy xác định. Khi được nung đến nhiệt độ nóng chảy trong điều kiện xác định, đồng chuyển từ trạng thái rắn sang trạng thái lỏng. Vì vậy, quá trình nóng chảy của đồng không diễn ra theo kiểu mềm dần trong một khoảng nhiệt độ như thủy tinh.
}
\end{bt}

% ID: L12C1B1-032-TL
% Mức: VD
%=== Câu 32 ===%
\dangbt{Ứng dụng và giải thích tính chất chất rắn}
\begin{bt}
Một cơ sở sản xuất cần chế tạo nhiều sản phẩm có hình dạng phức tạp. Người thợ dự định sử dụng hai vật liệu là thủy tinh và đồng.

\begin{enumerate}
\item Hãy phân loại thủy tinh và đồng thành chất rắn kết tinh hoặc chất rắn vô định hình.
\item So sánh đặc điểm nóng chảy của hai vật liệu.
\item Giải thích vì sao thủy tinh thích hợp cho phương pháp nung nóng rồi tạo hình trực tiếp.
\item Nếu muốn tạo sản phẩm bằng đồng có hình dạng phức tạp, hãy đề xuất một phương pháp gia công phù hợp.
\end{enumerate}

\loigiai{
\begin{enumerate}
\item Thủy tinh là chất rắn vô định hình; đồng là chất rắn kết tinh.

\item Thủy tinh không có nhiệt độ nóng chảy xác định và thường mềm dần trong một khoảng nhiệt độ khi nung nóng. Đồng có nhiệt độ nóng chảy xác định ở điều kiện xác định.

\item Do thủy tinh mềm dần khi nung nóng nên trong trạng thái mềm thích hợp, người thợ có thể uốn, nắn và tạo hình sản phẩm.

\item Có thể nung chảy đồng rồi đổ vào khuôn có hình dạng mong muốn. Ngoài ra, tùy sản phẩm, có thể sử dụng các phương pháp gia công cơ học như rèn, dập hoặc cán.
\end{enumerate}
}
\end{bt}
